% Proposed revision of the supplied experimental section.
% Keep the existing figure and section labels when integrating this file.
% Add gd-sgd-references.bib to the manuscript bibliography, or map the three
% citation keys below to entries already present in the main bibliography.
% AUTHOR CHECK: The supplied files contain the batch-size protocol but no
% results paragraph or underlying plot/data. The final subsection therefore
% describes the verified design and its interpretation without inventing a
% model-specific trend. Insert the observed results at the marked location.

\section{Gradient Reversal and Progress in GD and SGD}
\label{sec:gd-experiments}

Full-batch GD can continue reducing the training objective while successive
gradients become nearly antiparallel. In the three controlled problems
examined here, the largest tested learning rate achieves the lowest final
training loss within its model's sweep, although net displacement is only
2.6--5.4\% of the traveled parameter distance. We examine this coexistence
through consecutive alignment, directional curvature, cumulative motion,
and correlations across temporal lags. We then study how batch size changes
the alignment observed during SGD continuations.

Oscillatory descent at the edge of stability is established in prior
work~\cite{cohen2021edge,damian2023selfstabilization}; optimization trajectories
have also been studied through their directionality at different
timescales~\cite{singh2025directionality}. Our experiments connect the
gradient geometry of these oscillations to the distinction between local
reversal, accumulated parameter displacement, and training-loss reduction.

\paragraph{Controlled design.}
We train a tanh MLP on 5,000 CIFAR-10 images, ResNet-32 on all 50,000
CIFAR-10 training images, and a word-level Transformer on WikiText-2.
Each model minimizes its fixed training objective for 3,000 full-batch
updates. The constant learning rates are $\{0.004,0.02,0.04\}$ for the MLP,
$\{0.025,0.05,0.1\}$ for ResNet, and $\{0.05,0.1,0.2\}$ for the Transformer.
Within each model, the three rates share one frozen initialization and
data preparation. These comparisons isolate the effect of learning rate
within each controlled run; the batch-size study below uses three
initialization blocks. Appendix~\ref{appendix:gd-raw-alignment} specifies
the objectives and measurement conventions.

\subsection{The Onset of Sustained Gradient Reversal}

Figures~\ref{fig:gd-loss} and~\ref{fig:gd-alignment-displacement} show the
transition from positively aligned updates to oscillatory training.
The lowest MLP and Transformer rates produce smooth loss curves,
positive consecutive cosine, and $\eta\widehat\kappa_t$ near zero.
The other seven conditions develop nearly opposing consecutive gradients
while the loss continues to decrease over longer timescales.
ResNet exhibits this behavior at every tested rate, with an earlier
transition at larger learning rates.

The directional secant statistic expresses the same transition in
curvature terms. By Eq.~\eqref{eq:kappa-estimate}, nearly opposite gradients
with comparable consecutive norms give $\eta\widehat\kappa_t$ near two.
Alignment, adjacent norm ratio, and secant curvature therefore provide
linked descriptions of the observed geometry. In smooth full-batch GD,
the statistic measures directional curvature averaged along the preceding
step. One marks the sign change of consecutive alignment, and two is the
loss-change threshold in the quadratic approximation. The displayed
trailing means summarize regimes; the threshold interpretation applies
to the unsmoothed stepwise quantities.

The Hessian checkpoints in Fig.~\ref{fig:gd-loss} provide spectral context.
Their interpretation differs from that of the secant statistic in three
respects: they measure maximal rather than gradient-direction curvature,
they are pointwise rather than averaged along an update, and the ResNet
and Transformer measurements use fixed subset objectives. These differences
must be considered jointly when comparing the curves. The full-batch
gradient statistic remains available throughout each trajectory without
additional gradient evaluations or Hessian-vector products.

\subsection{Local Reversal Coexists with Substantial Loss Reduction}

Figure~\ref{fig:gd-alignment-displacement} relates consecutive reversal
to accumulated motion. In the MLP and ResNet, the transition is accompanied
by increasing gradient norms. Larger steps then contribute substantial
path length while opposing a component of the preceding update.
The normalized displacement $R_t$ decreases as traveled distance grows
faster than displacement from initialization.

After 3,000 updates at the highest tested rate, net displacement is
approximately 2.6\% of path length for the MLP, 3.7\% for ResNet, and
5.4\% for the Transformer. Yet each of these runs has the lowest final
training loss among its model's three tested rates. Thus, pronounced
local reversal and a small displacement-to-path-length ratio coexist
with substantial objective reduction in these experiments.

Here, ``destructive'' describes opposing parameter displacements.
The ratio $R_t$ measures the geometry of the traveled path; it does not
measure the fraction of learning retained. A small ratio can also arise
from changes in direction without opposing steps, so the evidence for
local reversal comes from the consecutive cosines. Likewise, the sign
of consecutive alignment alone does not determine the sum of interactions
across all temporal lags in Eq.~\eqref{eq:magnitude-interaction}.

\subsection{Reversal Has Structure across Temporal Lags}

We next examine whether negative consecutive alignment is part of a
recurring temporal pattern. Figure~\ref{fig:gd-backward} fixes the gradient
after 2,256 updates and compares it with the preceding 256 gradients,
using the highest rate for each model. The short lags alternate between
strongly negative odd-lag and strongly positive even-lag alignment.
The early envelope contracts more rapidly for the MLP and ResNet than
for the Transformer. Alignment can strengthen again at longer lags,
so the envelope is not monotone over the full window.

These measurements identify structured alternation at the selected
endpoint in each run. To relate this structure to displacement, note
that the exact motion over two consecutive updates is
\begin{equation}
  \theta_t-\theta_{t-2}=-\eta(g_{t-1}+g_t).
  \label{eq:gd-two-step-displacement}
\end{equation}
Opposing components cancel within this sum, while the remaining
component determines the two-step displacement. Its magnitude depends
on both alignment and the two gradient norms. The lag correlations
establish recurring reversal; the loss measurements separately
establish that the trajectory continues reducing the objective.
% AUTHOR: A paired-displacement analysis of the existing raw logs can
% quantify what survives each reversal and how those displacements add up.
% Do not insert its conclusion until the calculation has been performed.

\subsection{Observed Alignment under Mini-Batch Sampling}
\label{sec:sgd-noise-experiments}

The preceding full-batch runs isolate gradient geometry from sampling
noise. Figure~\ref{fig:sgd-moving-geometry} examines how the observed
alignment changes in 200-update continuations at different batch sizes,
started either at initialization or after 2,000 full-batch GD updates.
Each continuation contributes 199 adjacent sampled-gradient pairs.
Repetitions are averaged within each initialization, then the three
initialization blocks receive equal weight; shading shows their range.
The full-batch continuations provide a deterministic reference.

% AUTHOR: Specify the learning rate(s) used for these continuations, the
% learning rate used to produce each 2,000-update checkpoint, and the
% number of repetitions within each initialization. Neither supplied file
% provides these details. Do not infer them from the lag experiment.
% AUTHOR: Insert the model-specific result here, verified from the actual
% figure or data: describe which signs persist at large batch sizes,
% whether smaller batches move either sign toward zero, and any exception.
% Report seed variation when it changes the conclusion. Do not state that
% noise alone causes the observed trend in these moving trajectories.

The sampling unit is held consistent with the defined objective.
ResNet mini-batches contain whole fixed 1,000-image normalization groups;
Transformer batch sizes count prepared sequences. This distinction
matters when relating batch size to gradient sampling variance.

These continuations measure sampled-gradient alignment along evolving
trajectories. Changing batch size affects both the observed gradients
and the subsequent parameter states. Accordingly, a change in sampled
cosine can reflect a change in the visibility of the full-gradient
geometry, a change in the geometry itself, or both. This is the distinction
motivated by the noise decomposition in
Section~\ref{subsec:sampling-noise}; the present continuations characterize
the combined effect during SGD.
% AUTHOR: A complementary control can independently resample gradients
% at two frozen endpoints with known positive or negative full-gradient
% alignment. This holds the geometry fixed while testing sampling effects.
